\section{System Architecture}
\label{sec:architecture}

MemSt adopts a layered architecture (Figure \ref{fig:architecture}) with clear separation between interfaces, language bindings, and core storage components. This section describes the key architectural components.

\begin{figure*}[t]
    \centering
    \begin{tikzpicture}[
        node distance=1.5cm,
        box/.style={rectangle, draw, rounded corners, minimum width=3cm, minimum height=0.8cm, align=center, fill=lightblue!20},
        arrow/.style={->, >=stealth, thick}
    ]
        % Top layer
        \node[box, fill=memstblue!30] (ui) {Web UI / React};
        \node[box, fill=memstblue!30, right=of ui] (api) {REST API / FastAPI};
        \node[box, fill=memstblue!30, right=of api] (mcp) {MCP Adapter};
        \node[box, fill=memstblue!30, right=of mcp] (cli) {CLI / memst};
        
        % Middle layer
        \node[box, below=of ui, xshift=1.5cm] (python) {Python Bindings (PyO3)};
        \node[box, right=of python, xshift=1.5cm] (rust) {Rust Core (memst-core)};
        
        % Bottom layers
        \node[box, below=of python, fill=green!20] (search) {Hybrid Search\\BM25 + HNSW};
        \node[box, right=of search, fill=green!20] (repo) {MemRepo\\Git-like Store};
        \node[box, right=of repo, fill=green!20] (kg) {KG Extraction v2\\SQLite/DuckDB};
        \node[box, right=of kg, fill=green!20] (sleep) {Sleep-Time\\Consolidation};
        
        % Arrows
        \draw[arrow] (ui) -- (python);
        \draw[arrow] (api) -- (python);
        \draw[arrow] (mcp) -- (python);
        \draw[arrow] (cli) -- (rust);
        \draw[arrow] (python) -- (rust);
        \draw[arrow] (rust) -- (search);
        \draw[arrow] (rust) -- (repo);
        \draw[arrow] (rust) -- (kg);
        \draw[arrow] (rust) -- (sleep);
    \end{tikzpicture}
    \caption{MemSt architecture overview showing the multi-layer design with interfaces (top), language bindings (middle), and core storage/processing components (bottom).}
    \label{fig:architecture}
\end{figure*}

\subsection{Git-Like Memory Repository (MemRepo)}

At the foundation of MemSt is \textbf{MemRepo}, a content-addressable storage system inspired by Git's object model. Unlike traditional databases that use location-based addressing, MemRepo stores all data as immutable objects referenced by their Blake3 hash.

\subsubsection{Object Model}

MemRepo implements four fundamental object types:

\begin{itemize}
    \item \textbf{Blob}: Stores arbitrary binary content (messages, memories, embeddings). Each blob is addressed by $\text{Blake3}(\text{content})$.
    
    \item \textbf{Tree}: Represents hierarchical directory structures, mapping names to blob hashes and subtree hashes. Enables organized storage of session data.
    
    \item \textbf{Commit}: Captures a snapshot of the memory state at a point in time, containing a tree hash, parent commit hash(es), author/timestamp metadata, and a commit message.
    
    \item \textbf{Tag}: Named references to specific commits, enabling semantic labeling of memory milestones (e.g., ``session-start'', ``checkpoint-before-merge'').
\end{itemize}

This object model provides several advantages: (1) \textit{deduplication}---identical content stored once regardless of location; (2) \textit{immutability}---historical states are permanently preserved; (3) \textit{integrity}---any corruption is immediately detectable via hash mismatch.

\subsubsection{Write-Ahead Log (WAL)}

For durability and crash recovery, MemRepo implements a write-ahead log:

\begin{enumerate}
    \item All mutations are first appended to the WAL as journal entries
    \item Entries include operation type, affected objects, and rollback information
    \item Periodic checkpoints flush committed state to the object store
    \item On recovery, replay WAL entries from last checkpoint
\end{enumerate}

The WAL ensures ACID semantics: operations are atomic (all-or-nothing), consistent (integrity constraints enforced), isolated (via versioning), and durable (WAL persistence).

\subsubsection{Branching and Merging}

MemRepo supports Git-style branching for memory evolution:

\begin{itemize}
    \item \textbf{Branch}: A named reference to a commit, allowing parallel memory evolution
    \item \textbf{Checkout}: Switch the active branch to a different memory state
    \item \textbf{Merge}: Combine changes from two branches with three-way merge semantics
\end{itemize}

Branching enables powerful workflows: experimental memory configurations can be tested in isolation, agent behaviors can be versioned and rolled back, and multi-agent scenarios can use per-agent branches.

\subsection{Four-Tier Memory Hierarchy}

MemSt organizes memories into four tiers based on access frequency and importance:

\begin{table}[h]
\centering
\caption{Memory Tier Characteristics}
\begin{tabular}{@{}lllll@{}}
\toprule
\textbf{Tier} & \textbf{Access} & \textbf{Capacity} & \textbf{Promotion} & \textbf{Eviction} \\
\midrule
Working & Fastest & Limited & Manual & Age/importance \\
Short-term & Fast & Moderate & Access count & TTL \\
Long-term & Slower & Large & Importance score & Archival \\
Archival & Slowest & Unlimited & N/A & Manual \\
\bottomrule
\end{tabular}
\label{tab:memory-tiers}
\end{table}

\subsubsection{Working Memory}

Working memory contains the active conversation context and pinned facts. It is the only tier directly injected into the LLM context window. Working memory has strict token budget enforcement:

\begin{equation}
\sum_{m \in \text{Working}} \text{tokens}(m) \leq B_{\text{working}}
\end{equation}

where $B_{\text{working}}$ is configurable (default: 2048 tokens). When the budget is exceeded, memories are promoted to Short-term based on recency and importance scores.

\subsubsection{Short-Term Memory}

Short-term memory contains recent facts and user preferences that may be needed soon but are not immediately relevant. It is stored in compressed binary format with BM25 indexing for fast retrieval.

Promotion to Short-term occurs when:
\begin{itemize}
    \item Working memory exceeds its token budget
    \item A memory has been accessed more than $A_{\min}$ times (default: 3)
\end{itemize}

Demotion to Long-term occurs when a memory's time-to-live (TTL) expires (default: 7 days).

\subsubsection{Long-Term Memory}

Long-term memory stores important knowledge and historical data. Memories reach Long-term based on:

\begin{equation}
I(m) = \alpha \cdot \text{access\_freq}(m) + \beta \cdot \text{user\_rating}(m) + \gamma \cdot \text{semantic\_density}(m)
\end{equation}

where $I(m)$ is the importance score and $\alpha, \beta, \gamma$ are configurable weights. Memories with $I(m) > \theta_{\text{long}}$ are promoted to Long-term.

\subsubsection{Archival Memory}

Archival memory contains rarely accessed historical data. It is stored in highly compressed format and may be offloaded to slower storage (e.g., network-attached storage, S3). Archival memories retain full semantic indexing for retrieval but with higher latency.

\subsection{Context Assembly with Token Budget}

A unique feature of MemSt is \textbf{token-budget-aware context assembly}. When building the LLM prompt, MemSt does not simply retrieve the top-$k$ memories; it solves a constrained optimization problem:

\begin{align}
\text{maximize} \quad & \sum_{m \in S} R(m, q) \cdot I(m) \\
\text{subject to} \quad & \sum_{m \in S} \text{tokens}(m) \leq B_{\text{context}} - B_{\text{reserved}} \\
& |S| \leq N_{\max}
\end{align}

where $R(m, q)$ is the relevance of memory $m$ to query $q$, $I(m)$ is importance, $B_{\text{context}}$ is the model's context window, $B_{\text{reserved}}$ is tokens reserved for the response, and $N_{\max}$ is a maximum memory count for diversity.

This optimization ensures that the context window is filled with the most valuable information possible, rather than truncating an oversized retrieval set.

\subsection{Hybrid Search}

MemSt implements a three-backend search system:

\subsubsection{Native BM25}

A pure Rust implementation of the BM25 ranking function:

\begin{equation}
\text{BM25}(q, d) = \sum_{t \in q} \text{IDF}(t) \cdot \frac{f(t, d) \cdot (k_1 + 1)}{f(t, d) + k_1 \cdot (1 - b + b \cdot \frac{|d|}{\text{avgdl}})}
\end{equation}

with default parameters $k_1 = 1.2$, $b = 0.75$. This backend requires no external dependencies and is suitable for edge deployment.

\subsubsection{HNSW Vector Search}

For semantic similarity, MemSt uses Hierarchical Navigable Small World (HNSW) graphs:

\begin{itemize}
    \item Embeddings are computed via configurable API (OpenAI, local models)
    \item HNSW index enables approximate nearest neighbor search in $O(\log N)$
    \item Supports cosine similarity, Euclidean distance, and dot product
\end{itemize}

\subsubsection{Reciprocal Rank Fusion (RRF)}

To combine keyword and semantic signals, MemSt uses RRF:

\begin{equation}
\text{RRF}(d) = \sum_{r \in R} \frac{1}{k + r_r(d)}
\end{equation}

where $R$ is the set of rankers (BM25, HNSW), $r_r(d)$ is document $d$'s rank in ranker $r$, and $k = 60$ is a constant to dampen the impact of high rankings. This fusion approach is robust and requires no training.

\subsection{Multi-Agent Worktrees}

Inspired by Git worktrees, MemSt provides memory isolation for concurrent agents:

\begin{itemize}
    \item Each agent gets its own \textbf{worktree}---a separate working directory backed by the same MemRepo
    \item Worktrees can be created, switched, and removed without affecting other agents
    \item Changes in one worktree are invisible to others until explicitly merged
    \item The main worktree serves as a shared baseline
\end{itemize}

This architecture enables safe multi-agent deployment: agent A can experiment with memory modifications while agents B and C operate on stable baselines, with explicit merge gates controlling integration.
